\chapter{Objet, problématique et questions}\label{ch:objet}

\section{Les inscriptions datées comme documents transmis}

Une épitaphe ne livre ni la totalité d’une existence ni un relevé immédiat des croyances de ceux qui l’ont commandée. Elle agence un nom, une relation, une formule de commémoration et, parfois, une date. Ce qui en subsiste dépend du support, des pertes, des lectures et des opérations éditoriales successives. L’enquête doit donc porter à la fois sur les pratiques funéraires que les textes rendent accessibles et sur les médiations qui conditionnent cet accès.

Le volume rassemble \st{nA} notices relatives à des inscriptions classées comme chrétiennes et romaines, datées des années 340 à 403 dans les états documentaires étudiés. La distinction entre notice, texte, monument et personne gouverne toute l’analyse. Le catalogue conserve ses \st{nA} identifiants, afin que chaque renvoi reste stable. Les ensembles analytiques dépendent, eux, de décisions explicites sur les doublons, sur les fragments d’un même monument, sur la période et sur le genre des documents (chapitre~\ref{ch:corpus}). Aucun de ces ensembles ne constitue un recensement de la population chrétienne de Rome.

\section{Une question documentaire et une question historique}

La question documentaire porte sur les conditions de la preuve. Que permet d’affirmer un texte transmis par une base de données, lorsque son édition imprimée, sa photographie ou son support n’ont pas été contrôlés ? Que change la confrontation de deux bases indépendantes l’une de l’autre, mais qui dépendent toutes deux de la même édition ? La question historique porte sur les usages de la mémoire. Comment ces documents identifient-ils les défunts ? Comment représentent-ils leurs relations ? Comment inscrivent-ils la mort dans un langage chrétien ?

L’articulation de ces deux questions interdit deux raccourcis. Une restitution moderne ne peut devenir silencieusement une parole antique conservée. Inversement, l’absence d’un mot dans un texte bref ou lacunaire ne démontre pas l’absence de la croyance correspondante. Les désignations divines, les qualifications personnelles et le vocabulaire du repos sont donc examinés comme des formes d’énonciation. Ils ne sont pas convertis en mesure globale de l’intensité religieuse.

\section{Questions}

Le livre répond à cinq questions.

\begin{enumerate}[label=(\arabic*),leftmargin=2em]
\item \emph{Quel est le corpus ?} Combien de textes distincts les \st{nA} notices représentent-elles, et selon quels critères d’identité documentaire ? (chapitres~\ref{ch:corpus} et \ref{ch:transmission}).
\item \emph{Qu’est-ce qui date chaque texte ?} L’année repose-t-elle sur un nom consulaire lu, sur une lecture partielle, sur une restitution de l’éditeur ou sur la seule base de données ? Quel événement la date désigne-t-elle ? (partie~III).
\item \emph{Où et sur quoi ces textes ont-ils été gravés ?} Que permettent d’en dire les rattachements topographiques, les lieux de conservation et les descriptions matérielles de l’EDB ? (partie~IV).
\item \emph{Comment ces textes parlent-ils ?} Quelles sont la fréquence et l’évolution des formules funéraires, la langue des lapicides et les formes de la relation entre défunts et commémorants ? (parties~V et VI).
\item \emph{Que disent-ils de Dieu, du Christ et de la mort ?} Quel est le poids des mots gravés, des signes figurés et des restitutions dans ce qui est habituellement décrit comme le langage religieux des épitaphes ? (partie~VII).
\end{enumerate}

\section{Organisation de la démonstration}

L’enquête commence par la composition du corpus et par sa transmission numérique (partie~II). Elle distingue ensuite le repérage d’une année, le statut textuel de l’indication qui la fonde et l’événement daté (partie~III). L’analyse topographique et matérielle repose sur les données de l’EDB et reste limitée à leur résolution effective (partie~IV). Le formulaire, la langue, les noms, les relations familiales et les qualifications religieuses sont étudiés à partir de règles de relevé annoncées (parties~V à VII). La partie~VIII évalue la sensibilité des résultats à ces règles, puis les discute ; elle ne transforme pas un corpus documentaire sélectionné en échantillon aléatoire.

Les dossiers commentés et le catalogue raisonné (partie~IX) permettent de revenir des propositions générales aux passages qui les soutiennent. Les concordances et les index (partie~X) ont une fonction de repérage. Leur exhaustivité mécanique ne doit pas être confondue avec une validation philologique de chaque entrée.

\section{Statut du texte}

Les textes reproduits sont des lectures secondaires : ils sont donnés d’après l’EDCS, confrontés à l’EDB et, pour deux notices, à des reproductions universitaires (\nnums{38} et 68). Ni les volumes imprimés des \emph{Inscriptiones christianae urbis Romae}, ni les monuments n’ont été examinés. Ce statut est déclaré dans chaque notice. Il définit la portée du livre : une étude critique de la transmission d’un corpus daté, accompagnée d’analyses chronologiques, formulaires et historiques, et non une nouvelle édition critique des inscriptions. Les traductions sont des traductions de travail, à relire par un latiniste et un helléniste spécialistes.
